\section{Einleitung}
\label{sec:Einleitung}

\subsection{Vorwort}
\label{sub:Vorwort}
fghjfghjfghj TAI\footnote{Technischer Assistent für Informatik.} fghjfghj im Kapitel~\ref{sub:Projektumfeld} fghjfghj

\subsection{Kurzporträt der Ausbildungsstätte}
\label{sub:Kurzporträt der Ausbildungsstätte}
dfghdfghdfgh \Abkuerzung{UI} dfghdfgh

\begingroup
\setlength\intextsep{0pt}

\subsection{Projektumfeld}
\label{sub:Projektumfeld}
dfghdfghdfgh Non-SAP Technologien d dfghdfghdfgh Non-SAP Technologien d
dfghdfghdfgh Non-SAP Technologien d dfghdfghdfgh Non-SAP Technologien d
dfghdfghdfgh Non-SAP Technologien d dfghdfghdfgh Non-SAP Technologien d
dfghdfghdfgh Non-SAP Technologien d dfghdfghdfgh Non-SAP Technologien d

\begin{wrapfigure}[15]{r}[25pt]{0\textwidth}
    \includegraphicsKeepAspectRatio{./picture/systemlandschaft.pdf}{0.6}
    \caption{Systemlandschaft}
    \label{app:Systemlandschaft}
\end{wrapfigure}
sddfghdfgh \zB dfghdfgh (Siehe \Abbildung{app:Systemlandschaft})

dfghdfghdfgh Non-SAP Technologien d dfghdfghdfgh Non-SAP Technologien d
dfghdfghdfgh Non-SAP Technologien d dfghdfghdfgh Non-SAP Technologien d
dfghdfghdfgh Non-SAP Technologien d dfghdfghdfgh Non-SAP Technologien d
dfghdfghdfgh Non-SAP Technologien d dfghdfghdfgh Non-SAP Technologien d
dfghdfghdfgh Non-SAP Technologien d dfghdfghdfgh Non-SAP Technologien d
dfghdfghdfgh Non-SAP Technologien d dfghdfghdfgh Non-SAP Technologien d
dfghdfghdfgh Non-SAP Technologien d dfghdfghdfgh Non-SAP Technologien d
dfghdfghdfgh Non-SAP Technologien d dfghdfghdfgh Non-SAP Technologien d
dfghdfghdfgh Non-SAP Technologien d dfghdfghdfgh Non-SAP Technologien d
dfghdfghdfgh Non-SAP Technologien d dfghdfghdfgh Non-SAP Technologien d

\endgroup

\subsection{Motivation und Zielsetzung}
\label{sub:Motivation und Zielsetzung}
dfghdfgh

\subsection{Erklärung zum Dokument}
\label{sub:Erklärung zum Dokument}
An dieser Stelle wird die Gestaltung des Dokuments erläutert.
\begin{itemize}[noitemsep]
    \item Ein {[\textit{kursiver Text}]} weist auf eine Abkürzung hin, die im Abkürzungsverzeichnis zu finden ist.
    \item Ein {[\textbf{fettgedruckter Text}]}, mit Ausnahme von Überschriften, weist auf einen Fachbegriff hin, zu dem im Dokument
    eine Erklärung vorliegt. Ausgenommen sind Erklärungen in Fußnoten.
    \item Ein {[\texttt{Text in Schreibmaschinenschrift}]} weist auf eine Klasse oder Methode hin.
    \item Ein {[\textsf{Text in serifloser Schrift}]} weist auf sonstige Programmcode-Textbausteine hin.
    \item Liegt das Dokument in elektronischer Form vor, sind Abkürzungen und Quellenangaben zum jeweiligen Verzeichniseintrag
    verlinkt.
\end{itemize}

\subsection{Aufbau des Dokuments}
\label{sub:Aufbau des Dokuments}
dfghdfgh

